\DocumentMetadata{
  pdfversion=2.0,
  pdfstandard=ua-2,
  lang=en-US,
  tagging=on,
  tagging-setup={math/setup=mathml-SE,tabsorder=structure}
}
\documentclass[11pt,letterpaper]{article}
\usepackage[margin=0.68in,includefoot]{geometry}
\usepackage{newcomputermodern}
\usepackage{amsmath,amscd,mathtools}
\usepackage{tikz,adjustbox}
\providecommand{\square}{\mdlgwhtsquare}
\usepackage[hidelinks]{hyperref}
\hypersetup{
  pdftitle={Algebra PhD Exam, May 1993},
  pdfauthor={University of Florida Department of Mathematics},
  pdfsubject={Graduate mathematics examination},
  pdfdisplaydoctitle=true,
  bookmarksopen=true
}
\setcounter{secnumdepth}{0}
\setlength{\parindent}{0pt}
\setlength{\parskip}{0.28em}
\setlength{\emergencystretch}{3em}
\setlength{\leftmargini}{2.15em}
\setlength{\leftmarginii}{2.65em}
\setlength{\leftmarginiii}{3em}
\pagestyle{empty}
\raggedbottom
\allowdisplaybreaks
\NewTaggingSocketPlug{block/list/label}{examproblem}{%
  \tagstructbegin{tag=Lbl}\tagstructend%
  \tagstructbegin{tag=\LBody}%
  \tagstructbegin{tag=\ProblemHeadingTag,title-o={\ProblemHeadingText}}%
  \tagmcbegin{tag=\ProblemHeadingTag,actualtext={\ProblemHeadingText}}%
  #2%
  \tagmcend%
  \tagstructend%
}
\newenvironment{examproblems}[2]{%
  \def\ProblemHeadingTag{#2}%
  \AssignTaggingSocketPlug{block/list/label}{examproblem}%
  \begin{enumerate}%
  \setcounter{enumi}{#1}%
  \renewcommand{\labelenumi}{\textbf{\theenumi.}}%
  \setlength{\topsep}{0.35em}%
  \setlength{\partopsep}{0pt}%
  \setlength{\itemsep}{0.48em}%
  \setlength{\parsep}{0.18em}%
}{\end{enumerate}}
\newenvironment{examsubparts}{%
  \AssignTaggingSocketPlug{block/list/label}{default}%
  \begin{enumerate}%
  \setlength{\topsep}{0.18em}%
  \setlength{\partopsep}{0pt}%
  \setlength{\itemsep}{0.16em}%
  \setlength{\parsep}{0.1em}%
}{\end{enumerate}}
\newenvironment{examinstructions}{%
  \par\begingroup\itshape
}{%
  \par\endgroup\vspace{0.18em}
}
\makeatletter
\renewcommand\subsection{\@startsection{subsection}{2}{0pt}%
  {-1.25ex plus -.25ex minus -.1ex}%
  {0.55ex plus .1ex}%
  {\normalfont\large\bfseries}}
\renewcommand{\@maketitle}{%
  \newpage\null\vskip -1em%
  {\footnotesize\bfseries
    \href{https://www.ufl.edu/}{University of Florida}, \href{https://math.ufl.edu/}{Department of Mathematics}\par}%
  \vskip -0.5em%
  \tagstructbegin{tag=H1,title={Algebra PhD Exam, May 1993}}%
  \tagpdfparaOff%
  \tagmcbegin{tag=H1}%
    {\LARGE\bfseries\href{https://gma.math.ufl.edu/exams/algebra-phd/}{Algebra PhD Exam}, May 1993\par}%
  \tagmcend%
  \tagpdfparaOn%
  \tagstructend%
  \vskip 0.25em%
  \tagmcbegin{artifact}%
    \hrule height 1pt%
  \tagmcend%
  \vskip 1em%
}
\makeatother
\title{Algebra PhD Exam, May 1993}
\author{}
\date{}
\begin{document}
\maketitle
\thispagestyle{empty}
\begin{examinstructions}
Do seven of the following questions (the exam committee will NOT select your best seven answers if you answer more than seven). You may quote standard results (within reason) as long as you make it clear that you are doing so and you state them clearly.
\end{examinstructions}
\begin{examproblems}{0}{H2}
\def\ProblemHeadingText{Problem 1.}
\item
Define what is meant by a free group. Let~\(F\) be a free group on generators~\(X\) and let \(Y \subset X\). Prove that if~\(H\) is the smallest normal subgroup of~\(F\) generated by~\(Y\), then \(F/H\) is a free group.
\def\ProblemHeadingText{Problem 2.}
\item
Define what is meant by a nilpotent group. Show that if~\(G\) is a finite group, then~\(G\) is nilpotent if and only if every Sylow subgroup of~\(G\) is a normal subgroup of~\(G\).
\def\ProblemHeadingText{Problem 3.}
\item
State and prove Hilbert’s Theorem 90 about the elements of norm 1 in certain Galois extensions.
\def\ProblemHeadingText{Problem 4.}
\item
Let~\(F\) be a field of characteristic zero and let~\(M\) be a finite dimensional vector space over~\(F\). Let~\(b\) be a symmetric bilinear form on~\(M\). Prove that~\(M\) has a basis \(e_1,\ldots,e_n\) such that \(b(e_i,e_j)=0\) whenever \(i\ne j\).
\def\ProblemHeadingText{Problem 5.}
\item
Prove: let~\(S\) be an integral extension ring of the domain~\(R\), and let~\(P\) be a prime ideal of~\(R\). Then there is a prime ideal~\(Q\) of~\(S\) such that \(Q\cap R=P\).
\def\ProblemHeadingText{Problem 6.}
\item
Let~\(F\) be an algebraically closed field extension of a field~\(K\).
\begin{examsubparts}
\item[\textnormal{(i)}]
Prove that there is a one-to-one inclusion reversing correspondence between the set of affine~\(K\)-varieties in \(F^n\) and the set of radical ideals of \(K[x_1,\ldots,x_n]\).
\item[\textnormal{(ii)}]
If \(V_1\supset V_2\supset\cdots\) is a descending chain of~\(K\)-varieties in \(F^n\), then \(V_m=V_{m+1}=\cdots\) for some~\(m\).
\end{examsubparts}
\def\ProblemHeadingText{Problem 7.}
\item
Let~\(F\) be a field and let~\(A\) and~\(B\) be simple~\(F\)-algebras. Assume that~\(A\) is central (that is that \(Z(A)=F\cdot 1\)). Prove that \(A\otimes B\) is a simple~\(F\)-algebra.
\def\ProblemHeadingText{Problem 8.}
\item
Prove Wedderburn’s Theorem that every finite division ring is a field.
\def\ProblemHeadingText{Problem 9.}
\item
Let~\(A\) and~\(B\) be abelian groups. Prove
\begin{examsubparts}
\item[\textnormal{(i)}]
For each \(m>0\), \(A\otimes\mathbb{Z}_m\simeq A/mA\).
\item[\textnormal{(ii)}]
\(\mathbb{Z}_m\otimes\mathbb{Z}_n\simeq\mathbb{Z}_e\), where \(e=(m,n)\).
\item[\textnormal{(iii)}]
Describe \(A\otimes B\), when~\(A\) and~\(B\) are finitely generated.
\end{examsubparts}
\def\ProblemHeadingText{Problem 10.}
\item
Prove that no root of the polynomial \(x^5-6x-2\) can be obtained from rational numbers by using only the operations of addition, subtraction, multiplication, division and extraction of~\(n\)-th roots for various~\(n\).
\def\ProblemHeadingText{Problem 11.}
\item
Let~\(R\) be an associative ring.
\begin{examsubparts}
\item[\textnormal{(a)}]
What does it mean for the~\(R\)-module~\(P\) to be:
\begin{examsubparts}
\item[\textnormal{(i)}]
free?
\item[\textnormal{(ii)}]
projective?
\end{examsubparts}
\item[\textnormal{(b)}]
Are free modules always projective? Are projective modules always free? Justify your answers.
\item[\textnormal{(c)}]
Give an example of a noncommutative ring for which the projective modules and the free modules are the same.
\end{examsubparts}
\end{examproblems}
\end{document}
